\documentclass[14pt]{extarticle}
\usepackage{natbib}

\usepackage{xcolor}
\usepackage{amsmath}
\newcommand{\tuple}[1]{ \langle #1 \rangle }
%\usepackage{automata}
\usepackage{times}
\usepackage{ltablex}
\usepackage{tasks}
\usepackage{threeparttable}
\usepackage{booktabs}
\usepackage{xcolor} % For custom colors
\usepackage{tikz} % For styling enumerate numbers
\usepackage{tcolorbox} % For colored box styling
\usepackage{amsmath, amsfonts, amssymb, amsthm} % Math related
\usepackage{natbib}
\usepackage{fontspec}
\usepackage{luatexja}
\usepackage[mathscr]{euscript}

%%%%%% Template
\usepackage{hyperref}
\hypersetup{colorlinks=true,allcolors=blue}

\usepackage{vmargin}
\setpapersize{USletter}
\setmarginsrb{1.0in}{1.0in}{1.0in}{0.6in}{0pt}{0pt}{0pt}{0.4in}

% HOW TO USE THE ABOVE:
%\setmarginsrb{leftmargin}{topmargin}{rightmargin}{bottommargin}{headheight}{headsep}{footheight}{footskip}
%\raggedbottom
% paragraphs indent & skip:
\parindent  0.3cm
\parskip    -0.01cm

\usepackage{tikz}
\usetikzlibrary{backgrounds}

% hyphenation:
% \hyphenpenalty=10000 % no hyphen
% \exhyphenpenalty=10000 % no hyphen
\sloppy

% notes-style paragraph spacing and indentation:
\usepackage{parskip}
\setlength{\parindent}{0cm}

% let derivations break across pages
\allowdisplaybreaks

\newcommand{\orange}[1]{\textcolor{orange}{#1}}
\newcommand{\blue}[1]{\textcolor{blue}{#1}}
\newcommand{\red}[1]{\textcolor{red}{#1}}
\newcommand{\freq}[1]{{\bf \sf F}(#1)}
\newcommand{\datafreq}[2]{{{\bf \sf F}_{#1}(#2)}}

\def\qqquad{\quad\qquad}
\def\qqqquad{\qquad\qquad}

%%%%%%%%%%%%%%%%%%%%%%%%%%%%%%%%%%%%%%%%%%%%%%%%%%%%%%%%%%%%%%%%%%%%%%%%%%%%%%%%
%%%%%%%%%%%%%%%%%%%%%%%%%%%%%%%%%%%%%%%%%%%%%%%%%%%%%%%%%%%%%%%%%%%%%%%%%%%%%%%%

% fill-in-blank question style, found in https://tex.stackexchange.com/a/505089

\usepackage{ifthen}
\usepackage{tocloft}
\usepackage{exercise}
% \usepackage{xcolor}

% Set the Show Answers Boolean
\newboolean{showAns}
\setboolean{showAns}{false}
\newcommand{\showAns}{\setboolean{showAns}{true}}

% The length of the Answer line
\newlength{\answerlength}
\newcommand{\anslen}[1]{\settowidth{\answerlength}{#1}}

% ans command that indicates space for an answer or shows the answer in red
\newcommand{\ans}[1]{\settowidth{\answerlength}{\hspace{2ex}#1\hspace{2ex}}%
    \ifthenelse{\boolean{showAns}}%
        {\textcolor{red}{\underline{\hspace{2ex}#1\hspace{2ex}}}}%
        {\underline{\hspace{\answerlength}}}}%

\newcommand{\details}[1]{\settowidth{\answerlength}{#1}%
    \ifthenelse{\boolean{showAns}}%
        {\begin{tcolorbox}[colback=blue!5!white,colframe=blue!75!black,title=Solution] #1 \end{tcolorbox}}%
        {}}%

% Formatting how multiple choices Questions are formated.
\settasks{label=(\Alph*), label-width=30pt}


% Some commands for the Exercise Question package
\renewcommand{\QuestionNB}{\Large\protect\textcircled{\small\bfseries\arabic{Question}}\ }
\renewcommand{\ExerciseHeader}{} %no header
\renewcommand{\QuestionBefore}{3ex} %Space above each Q
\setlength{\QuestionIndent}{8pt} % Indent after Q number


% To create the list of answers with tocloft...
\newcommand{\listanswername}{Answers}
\newlistof[Question]{answer}{Answers}{\listanswername}

% Creates a TOC for Answers
\newcounter{prevQ}
\newcommand{\answer}[1]{\refstepcounter{answer}%
\ans{#1}%
\ifnum\theQuestion=\theprevQ%
        \addcontentsline{Answers}{answer}{\protect\numberline{}#1}% don't include the Q number
        \else%
        \addcontentsline{Answers}{answer}{\protect\numberline{\theQuestion}#1}%
        \setcounter{prevQ}{\value{Question}}%
        \fi%
        }%

% \hyphenpenalty=10000 % no hyphen
% \exhyphenpenalty=10000 % no hyphen
\sloppy              % hyphen

\newcommand{\HRule}{\rule{\linewidth}{0.5mm}}
\newcommand{\Hrule}{\rule{\linewidth}{0.3mm}}

%tocloft formatting listofanswers
\renewcommand{\cftAnswerstitlefont}{\bfseries\large}
\renewcommand{\cftanswerdotsep}{\cftnodots}
\cftpagenumbersoff{answer}
\addtolength{\cftanswernumwidth}{10pt}

\makeatletter% since there's an at-sign (@) in the command name
\renewcommand{\@maketitle}{%
  \parindent=0pt% don't indent paragraphs in the title block
  \centering
  {\Large \bfseries\textsc{\@title}} \\
  \vspace{5pt}
  {\large \textit{\@author}} \\
  \HRule \\
  \vspace{1em}
}
\makeatother% resets the meaning of the at-sign (@)


% --- Fonts (robust fallback) ---
\IfFontExistsTF{Crimson Pro Light}{
  \setmainfont{Crimson Pro Light}[
    ItalicFont={* Italic},
    BoldFont={Crimson Pro Medium},
    BoldItalicFont={Crimson Pro Medium Italic}]
  \setsansfont{Crimson Pro Light}[
    ItalicFont={* Italic},
    BoldFont={Crimson Pro Medium},
    BoldItalicFont={Crimson Pro Medium Italic}]
}{
  \setmainfont{HaranoAjiMincho}
  \setsansfont{HaranoAjiGothic}
}
\title{Final Exam}
\author{Macroeconomics II \\ Hui-Jun Chen}


% ------------------------------------------------------------
% Explanations for each option (only printed when \showAns is enabled)
% Usage: \ExplainChoices{A}{B}{C}{D}
% ------------------------------------------------------------
\newcommand{\ExplainChoices}[4]{%
\ifthenelse{\boolean{showAns}}{%
\vspace{0.3em}\noindent\textcolor{red}{\textbf{Explanations.}}
\begin{itemize}
\item[(A)] #1
\item[(B)] #2
\item[(C)] #3
\item[(D)] #4
\end{itemize}
}{}%
}
\begin{document}

\maketitle

% Uncomment to display answers + explanations
\showAns

\begin{Exercise}
\section*{Problem 1 (Questions 1--20): Measurement, IS--LM, AD--AS}

\Question \textbf{Units. If CPI increases from 180 to 189 over a year, inflation is closest to}
\answer{B}
\begin{tasks}(1)
\task $2\%$
\task $5\%$
\task $9\%$
\task $-5\%$
\end{tasks}
\ExplainChoices{Wrong: too small.}{Correct: $(189-180)/180=0.05$.}{Wrong: 9\% would be 180$\to$196.2.}{Wrong sign.}

\Question \textbf{Approximation. For small changes, $\log(P_t/P_{t-1})$ is approximately}
\answer{B}
\begin{tasks}(1)
\task $P_t-P_{t-1}$
\task $(P_t-P_{t-1})/P_{t-1}$
\task $(P_t-P_{t-1})/P_t$
\task $P_{t-1}/P_t$
\end{tasks}
\ExplainChoices{Wrong: not scale-free.}{Correct: first-order approximation.}{Wrong denominator.}{Wrong.}

\Question \textbf{Fisher (numbers). If $i=1.5\%$ and $\mathbb{E}\pi=2.0\%$, the ex-ante real rate is}
\answer{C}
\begin{tasks}(1)
\task $3.5\%$
\task $0.5\%$
\task $-0.5\%$
\task $2.0\%$
\end{tasks}
\ExplainChoices{Wrong: adds.}{Wrong sign.}{Correct: $1.5-2.0=-0.5$.}{Wrong: equals expected inflation.}

\Question \textbf{ZLB. If $i=0$ and expected inflation falls from $1\%$ to $-1\%$, the real rate changes by}
\answer{C}
\begin{tasks}(1)
\task $-2$ pp
\task $0$ pp
\task $+2$ pp
\task $+1$ pp
\end{tasks}
\ExplainChoices{Wrong sign.}{Wrong: $r=-\mathbb{E}\pi$ at ZLB.}{Correct: real rate rises by 2 percentage points.}{Wrong.}

\Question \textbf{Debt valuation view. The price level tends to rise when expected future primary surpluses}
\answer{B}
\begin{tasks}(1)
\task increase
\task fall
\task stay constant
\task equal current taxes
\end{tasks}
\ExplainChoices{Wrong: more backing raises real value.}{Correct: weaker backing raises $P$ to revalue nominal claims.}{Not necessarily.}{Not the key condition.}

\Question \textbf{Okun. If $u-u^n=-0.4x$ and $x=-3\%$, then $u-u^n$ is}
\answer{B}
\begin{tasks}(1)
\task $-1.2\%$
\task $+1.2\%$
\task $-0.75\%$
\task $+0.75\%$
\end{tasks}
\ExplainChoices{Wrong sign.}{Correct: $-0.4\times(-3)=+1.2$.}{Wrong.}{Wrong.}

\Question \textbf{Money market numerics. If $M/P=12$, $k=0.8$, $h=4$, and $Y=20$, then $i$ equals}
\answer{A}
\begin{tasks}(1)
\task $1$
\task $2$
\task $3$
\task $4$
\end{tasks}
\ExplainChoices{Correct: $i=(k/h)Y-(1/h)(M/P)=(0.2)20-0.25\cdot12=1$.}{Wrong.}{Wrong.}{Wrong.}

\Question \textbf{IS numerics. If $Y=200-50r$ and $r$ falls by 0.01, then $\Delta Y=$}
\answer{B}
\begin{tasks}(1)
\task $-0.5$
\task $+0.5$
\task $+5$
\task $-5$
\end{tasks}
\ExplainChoices{Wrong sign.}{Correct: $-50\times(-0.01)=+0.5$.}{Wrong magnitude.}{Wrong sign/magnitude.}

\Question \textbf{AD slope (IS--LM). With $M$ fixed, $P\uparrow$ leads to $Y\downarrow$ because}
\answer{B}
\begin{tasks}(1)
\task real balances rise
\task real balances fall and $i$ rises
\task IS shifts left mechanically
\task $Y^{FE}$ rises
\end{tasks}
\ExplainChoices{Wrong.}{Correct.}{Wrong.}{Wrong.}

\Question \textbf{AD--AS numerics. AD: $y=50-2p$, SRAS: $p=1+0.5(y-40)$. Equilibrium output $y$ is}
\answer{C}
\begin{tasks}(1)
\task $40$
\task $42$
\task $44$
\task $46$
\end{tasks}
\ExplainChoices{Wrong.}{Wrong.}{Correct: solve to get $y=44$.}{Wrong.}

\Question \textbf{Using the previous AD--AS system, the equilibrium price level $p$ is}
\answer{C}
\begin{tasks}(1)
\task $1$
\task $2$
\task $3$
\task $4$
\end{tasks}
\ExplainChoices{Wrong.}{Wrong.}{Correct: $p=0.5\cdot44-19=3$.}{Wrong.}

\Question \textbf{Diagnosis. A positive demand shock in AD--AS typically causes}
\answer{A}
\begin{tasks}(1)
\task $y\uparrow,\,p\uparrow$
\task $y\downarrow,\,p\uparrow$
\task $y\uparrow,\,p\downarrow$
\task $y\downarrow,\,p\downarrow$
\end{tasks}
\ExplainChoices{Correct.}{Supply shock pattern.}{Wrong.}{Wrong.}

\Question \textbf{Diagnosis. A negative supply shock in AD--AS typically causes}
\answer{B}
\begin{tasks}(1)
\task $y\uparrow,\,p\uparrow$
\task $y\downarrow,\,p\uparrow$
\task $y\uparrow,\,p\downarrow$
\task $y\downarrow,\,p\downarrow$
\end{tasks}
\ExplainChoices{Wrong.}{Correct (stagflation).}{Wrong.}{Wrong.}

\Question \textbf{Medium run. Output returns toward $y^{FE}$ mainly because}
\answer{B}
\begin{tasks}(1)
\task AD shifts back automatically
\task expected prices/wages adjust (SRAS shifts)
\task money supply is constant
\task productivity always rises
\end{tasks}
\ExplainChoices{Wrong.}{Correct.}{Wrong.}{Wrong.}

\Question \textbf{IS--LM: a cut in lump-sum taxes (holding $G$ fixed) is typically modeled as}
\answer{A}
\begin{tasks}(1)
\task IS shifts right
\task LM shifts right
\task IS shifts left
\task LM shifts left
\end{tasks}
\ExplainChoices{Correct: autonomous demand rises.}{Wrong.}{Wrong.}{Wrong.}

\Question \textbf{IS--LM: a decrease in nominal money $M$ (holding $P$ fixed) shifts}
\answer{B}
\begin{tasks}(1)
\task LM right/down
\task LM left/up
\task IS right
\task IS left
\end{tasks}
\ExplainChoices{Wrong.}{Correct.}{Wrong.}{Wrong.}

\Question \textbf{Liquidity trap logic. If the money-market-clearing formula implies $i^\ast<0$ and there is a ZLB, then the equilibrium $i$ is}
\answer{B}
\begin{tasks}(1)
\task $i^\ast$
\task $0$
\task $|i^\ast|$
\task undetermined always
\end{tasks}
\ExplainChoices{Wrong: violates ZLB.}{Correct.}{Wrong.}{Wrong.}

\Question \textbf{One-time jump. If $P$ rises by 6\% once and then stays constant, inflation is}
\answer{B}
\begin{tasks}(1)
\task 6\% forever
\task 6\% once then 0\%
\task 0\% forever
\task -6\% once then 0\%
\end{tasks}
\ExplainChoices{Wrong.}{Correct.}{Wrong.}{Wrong sign.}

\Question \textbf{Anchored expectations means}
\answer{A}
\begin{tasks}(1)
\task $\mathbb{E}\pi$ stays near target after temporary shocks
\task $\pi=0$ always
\task $x=0$ always
\task $i$ never changes
\end{tasks}
\ExplainChoices{Correct.}{Wrong.}{Wrong.}{Wrong.}

\Question \textbf{IS--MP--PC uses which three blocks?}
\answer{B}
\begin{tasks}(1)
\task IS, LM, Phillips
\task IS, policy rule, Phillips
\task AD, money demand, LRAS
\task budget constraint, money demand, Phillips
\end{tasks}
\ExplainChoices{Wrong: LM replaced by a rule.}{Correct.}{Wrong.}{Wrong.}

\vspace{0.6em}\Hrule\vspace{0.8em}
\section*{Problem 2 (Questions 21--40): IS--MP--PC, Taylor principle, NK}

\Question \textbf{In Whelan’s IS--MP derivation, the key term is $(\beta_\pi-1)$ because the real-rate gap contains}
\answer{A}
\begin{tasks}(1)
\task $i-\pi$
\task $i+\pi$
\task $\pi-i$
\task $i/\pi$
\end{tasks}
\ExplainChoices{Correct: $r=i-\pi$ (or $i-\pi^e$).}{Wrong.}{Wrong.}{Wrong.}

\Question \textbf{Taylor principle. $\beta_\pi>1$ ensures that when inflation rises, the real rate}
\answer{B}
\begin{tasks}(1)
\task falls
\task rises
\task is unchanged
\task must be zero
\end{tasks}
\ExplainChoices{Wrong.}{Correct.}{Wrong.}{Wrong.}

\Question \textbf{If $\theta=\dfrac{1}{1+\alpha\gamma(\beta_\pi-1)}$ and $\beta_\pi>1$, then}
\answer{C}
\begin{tasks}(1)
\task $\theta>1$
\task $\theta=1$
\task $0<\theta<1$
\task $\theta<0$
\end{tasks}
\ExplainChoices{Wrong.}{Wrong.}{Correct.}{Wrong.}

\Question \textbf{Numerics. If $\alpha=1.2$, $\gamma=0.4$, $\beta_\pi=2$, then $\theta$ is closest to}
\answer{A}
\begin{tasks}(1)
\task $0.68$
\task $0.32$
\task $1.68$
\task $0$
\end{tasks}
\ExplainChoices{Correct: denom $1+1.2\cdot0.4\cdot1=1.48$.}{Wrong.}{Wrong.}{Wrong.}

\Question \textbf{A tougher central bank corresponds to}
\answer{B}
\begin{tasks}(1)
\task lower $\beta_\pi$
\task higher $\beta_\pi$
\task higher $\gamma$
\task lower $\alpha$
\end{tasks}
\ExplainChoices{Wrong.}{Correct.}{Wrong: $\gamma$ is Phillips slope.}{Not the toughness definition.}

\Question \textbf{With a tougher CB, a cost-push shock tends to produce}
\answer{B}
\begin{tasks}(1)
\task larger inflation rise, smaller output fall
\task smaller inflation rise, larger output fall
\task smaller inflation rise, smaller output fall
\task no changes
\end{tasks}
\ExplainChoices{Wrong.}{Correct.}{Wrong.}{Wrong.}

\Question \textbf{Explosive region in the TaylorRule discussion: if $1-\frac{1}{\alpha\gamma}<\beta_\pi<1$, then}
\answer{B}
\begin{tasks}(1)
\task IS--MP slopes down and is stable
\task IS--MP slopes up and can be unstable under adaptive expectations
\task inflation is pinned at target
\task Phillips curve becomes vertical
\end{tasks}
\ExplainChoices{Wrong.}{Correct.}{Wrong.}{Wrong.}

\Question \textbf{Boundary. If $\alpha\gamma=0.5$, then $1-\frac{1}{\alpha\gamma}=$}
\answer{B}
\begin{tasks}(1)
\task $0.5$
\task $-1$
\task $1$
\task $2$
\end{tasks}
\ExplainChoices{Wrong.}{Correct: $1-2=-1$.}{Wrong.}{Wrong.}

\Question \textbf{Two-period NK IS (calculation). With $x_1=0$, $\sigma=1$, $r_0^n=1\%$, $i_0=4\%$, and $\mathbb{E}_0\pi_1=2\%$, $x_0$ is}
\answer{C}
\begin{tasks}(1)
\task $+1\%$
\task $0\%$
\task $-1\%$
\task $-3\%$
\end{tasks}
\ExplainChoices{Wrong sign.}{Wrong.}{Correct: $4-2-1=1$ so $x_0=-1\%$.}{Wrong.}

\Question \textbf{Forward-looking NKPC is}
\answer{B}
\begin{tasks}(1)
\task $\pi_0=\pi_{-1}+\kappa x_0+u_0$
\task $\pi_0=\beta\mathbb{E}_0\pi_1+\kappa x_0+u_0$
\task $\pi_0=\kappa x_0$ always
\task $\pi_0=\beta\pi_0+\kappa x_0$
\end{tasks}
\ExplainChoices{Backward-looking.}{Correct.}{Only if expectations and shocks are zero.}{Wrong.}

\Question \textbf{NKPC numerics. If $\beta=0.99$, $\mathbb{E}_0\pi_1=1\%$, $\kappa=0.3$, $x_0=2\%$, $u_0=0$, then $\pi_0$ is}
\answer{A}
\begin{tasks}(1)
\task $1.59\%$
\task $1.29\%$
\task $0.99\%$
\task $0.60\%$
\end{tasks}
\ExplainChoices{Correct: $0.99+0.6=1.59$.}{Wrong.}{Wrong: ignores gap term.}{Wrong: gap term only.}

\Question \textbf{Taylor rule numerics. If $i_0=\bar i+\phi_\pi\pi_0$ with $\bar i=1\%$, $\phi_\pi=1.5$, and $\pi_0=3\%$, then $i_0$ is}
\answer{C}
\begin{tasks}(1)
\task $3\%$
\task $4.5\%$
\task $5.5\%$
\task $6\%$
\end{tasks}
\ExplainChoices{Wrong.}{Wrong (missing intercept).}{Correct: $1+4.5=5.5$.}{Wrong.}

\Question \textbf{Two-period Fisher numerics. If $i_0=0\%$ and $\mathbb{E}_0\pi_1=2\%$, then $r_0$ is}
\answer{C}
\begin{tasks}(1)
\task $2\%$
\task $0\%$
\task $-2\%$
\task $-4\%$
\end{tasks}
\ExplainChoices{Wrong sign.}{Wrong.}{Correct.}{Wrong.}

\Question \textbf{Two-period IS sign. Holding $\mathbb{E}_0\pi_1$ fixed, a decrease in $i_0$ implies}
\answer{B}
\begin{tasks}(1)
\task $x_0$ decreases
\task $x_0$ increases
\task $x_0$ unchanged
\task $\pi_0$ falls one-for-one
\end{tasks}
\ExplainChoices{Wrong.}{Correct.}{Wrong.}{Wrong.}

\Question \textbf{In NKPC, increasing $u_0$ is best interpreted as}
\answer{B}
\begin{tasks}(1)
\task a right shift in IS
\task an upward shift in inflation for given $x_0$ and expectations
\task a fall in $r_0^n$
\task a decrease in $\kappa$
\end{tasks}
\ExplainChoices{Wrong.}{Correct.}{Wrong.}{Wrong.}

\Question \textbf{Two-period closure. If $x_1=0$ and $\pi_1=0$ are imposed, then at $t=0$}
\answer{A}
\begin{tasks}(1)
\task there are no expectations terms
\task the system is purely backward-looking
\task policy becomes irrelevant
\task the LM curve returns
\end{tasks}
\ExplainChoices{Correct: $\mathbb{E}_0x_1=\mathbb{E}_0\pi_1=0$.}{Wrong.}{Wrong.}{Wrong.}

\Question \textbf{Determinacy in basic NK. A standard sufficient condition is}
\answer{B}
\begin{tasks}(1)
\task $\phi_\pi<1$
\task $\phi_\pi>1$
\task $\kappa=0$
\task $\sigma=0$
\end{tasks}
\ExplainChoices{Wrong.}{Correct (Taylor principle).}{Wrong.}{Wrong.}

\Question \textbf{Two-period NK IS (calculation). Suppose $x_1=0$, $\sigma=1$, $r_0^n=1\%$, $i_0=3\%$, and $\mathbb{E}_0\pi_1=0\%$. Then $x_0$ is}
\answer{B}
\begin{tasks}(1)
\task $+2\%$
\task $-2\%$
\task $-1\%$
\task $0\%$
\end{tasks}
\ExplainChoices{Wrong sign.}{Correct: $3-0-1=2$ so $x_0=-2\%$.}{Wrong.}{Wrong.}

\Question \textbf{IS--MP (calculation). If $\alpha=1.5$ and $\beta_\pi=0.8$, then $-\alpha(\beta_\pi-1)$ equals}
\answer{B}
\begin{tasks}(1)
\task $-0.3$
\task $+0.3$
\task $-1.2$
\task $+1.2$
\end{tasks}
\ExplainChoices{Wrong sign.}{Correct: $-1.5(-0.2)=+0.3$.}{Wrong magnitude.}{Wrong magnitude.}

\Question \textbf{IS--MP slope (sign). If $\beta_\pi=0.6$ and $\alpha>0$, then IS--MP in $(y,\pi)$ space is}
\answer{B}
\begin{tasks}(1)
\task downward sloping
\task upward sloping
\task vertical
\task flat
\end{tasks}
\ExplainChoices{Wrong: downward requires $\beta_\pi>1$.}{Correct: $\beta_\pi<1$ implies positive slope.}{Wrong.}{Wrong.}

\end{Exercise}
\end{document}
